\begin{textsample}{NEG Dim 1 – human – Score -6.00 – mg/mg\_ef\_linguagens\_ed\_fis\_1\_p5.txt}  \label{ex:f1_neg_018}
A localização desse componente curricular na área de Linguagens também delimita uma mudança de paradigma e de base conceitual pois acolhe a gestualidade humana como algo mais amplo do que as relações neuromotoras consequentes dos processos bio-fisiológicos vinculados ao \textbf{movimento} humano.
Ela é compreendida, isto sim, como \textbf{forma} e meio de expressão e comunicação repleta de \textbf{sentidos} e \textbf{significados} \textbf{produzidos} e reproduzidos, criados e recriados nos espaços de mediação social e cultural.
Daí surge a delimitação do objeto de estudo da educação física escolar : as práticas \textbf{corporais} organizadas como Cultura \textbf{Corporal} de \textbf{Movimentos}.

% matched lemmas: corporal, forma, movimento, produzir, sentido, significado
\end{textsample}
